\section{Requisiti di Sistema}\label{requisiti-di-sistema}

I requisiti di sistema definiscono i vincoli architetturali,
infrastrutturali e operativi che BlueAgent deve garantire per supportare
i casi d'uso descritti. Di seguito vengono presentati, suddivisi nelle
rispettive categorie funzionali, i requisiti formali di progetto.

\subsection{Piattaforma Host e
Infrastruttura}\label{piattaforma-host-e-infrastruttura}

La tabella seguente riassume i requisiti relativi al sistema operativo
host e ai componenti infrastrutturali di base necessari all'esecuzione
di BlueAgent.

\begin{longtable}[]{@{}
  >{\raggedright\arraybackslash}p{(\linewidth - 4\tabcolsep) * \real{0.1531}}
  >{\raggedright\arraybackslash}p{(\linewidth - 4\tabcolsep) * \real{0.2245}}
  >{\raggedright\arraybackslash}p{(\linewidth - 4\tabcolsep) * \real{0.6224}}@{}}
\toprule\noalign{}
\begin{minipage}[b]{\linewidth}\raggedright
Codice
\end{minipage} & \begin{minipage}[b]{\linewidth}\raggedright
Requisito
\end{minipage} & \begin{minipage}[b]{\linewidth}\raggedright
Descrizione Operativa
\end{minipage} \\
\midrule\noalign{}
\endhead
\bottomrule\noalign{}
\endlastfoot
\textbf{\passthrough{\lstinline!SR-HST-01!}} & Sistema Operativo Host &
L'host deve eseguire un kernel Linux, con IP forwarding abilitato e
supporto completo per il sottosistema
\passthrough{\lstinline!nftables!}. \\
\textbf{\passthrough{\lstinline!SR-HST-02!}} & Storage Pool ZFS & L'host
deve disporre di uno storage pool ZFS
(\passthrough{\lstinline!blueagent-zfs!}) con supporto nativo a snapshot
Copy-on-Write (CoW). \\
\textbf{\passthrough{\lstinline!SR-HST-03!}} & Servizi SDN Host & L'host
deve eseguire i demoni Open vSwitch e OVN
(\passthrough{\lstinline!ovs-vswitchd!},
\passthrough{\lstinline!ovn-northd!},
\passthrough{\lstinline!ovn-controller!}) per il provisioning dinamico
degli switch logici effimeri. \\
\textbf{\passthrough{\lstinline!SR-HST-04!}} & Supporto WireGuard &
L'host deve fornire il modulo kernel WireGuard e i relativi strumenti di
configurazione. \\
\textbf{\passthrough{\lstinline!SR-HST-05!}} & Pipeline Distrobuilder &
L'host deve disporre di \passthrough{\lstinline!distrobuilder!} per la
compilazione automatizzata e riproducibile delle immagini dei
container. \\
\end{longtable}

\subsection{Motore di
Containerizzazione}\label{motore-di-containerizzazione}

I requisiti per il motore di containerizzazione stabiliscono i
meccanismi di interazione con Incus, l'isolamento dei carichi e la
gestione del ciclo di vita dei container di sistema.

\begin{longtable}[]{@{}
  >{\raggedright\arraybackslash}p{(\linewidth - 4\tabcolsep) * \real{0.1531}}
  >{\raggedright\arraybackslash}p{(\linewidth - 4\tabcolsep) * \real{0.2245}}
  >{\raggedright\arraybackslash}p{(\linewidth - 4\tabcolsep) * \real{0.6224}}@{}}
\toprule\noalign{}
\begin{minipage}[b]{\linewidth}\raggedright
Codice
\end{minipage} & \begin{minipage}[b]{\linewidth}\raggedright
Requisito
\end{minipage} & \begin{minipage}[b]{\linewidth}\raggedright
Descrizione Operativa
\end{minipage} \\
\midrule\noalign{}
\endhead
\bottomrule\noalign{}
\endlastfoot
\textbf{\passthrough{\lstinline!SR-ENG-01!}} & Astrazione del Motore &
Disaccoppiamento della logica di orchestrazione tramite l'interfaccia
astratta \passthrough{\lstinline!ContainerEngine!}. \\
\textbf{\passthrough{\lstinline!SR-ENG-02!}} & Client Asincrono UDS &
Comunicazione non bloccante con Incus via Unix Domain Socket
(\passthrough{\lstinline!/var/lib/incus/unix.socket!}) mediante client
HTTP asincrono. \\
\textbf{\passthrough{\lstinline!SR-ENG-03!}} & Limitazione Risorse
(Cgroups v2) & Applicazione di quote hardware per impedire la starvation
delle risorse dell'host. \\
\textbf{\passthrough{\lstinline!SR-ENG-04!}} & Clonazione Snapshot CoW &
Duplicazione istantanea delle istanze container basata su snapshot CoW a
livello di storage pool ZFS. \\
\end{longtable}

\subsection{Rete Ibrida e
Connettività}\label{rete-ibrida-e-connettivituxe0}

Questa sezione elenca i requisiti di rete necessari a garantire
prestazioni native in fase di gara e isolamento crittografico per ogni
squadra.

\begin{longtable}[]{@{}
  >{\raggedright\arraybackslash}p{(\linewidth - 4\tabcolsep) * \real{0.1531}}
  >{\raggedright\arraybackslash}p{(\linewidth - 4\tabcolsep) * \real{0.2245}}
  >{\raggedright\arraybackslash}p{(\linewidth - 4\tabcolsep) * \real{0.6224}}@{}}
\toprule\noalign{}
\begin{minipage}[b]{\linewidth}\raggedright
Codice
\end{minipage} & \begin{minipage}[b]{\linewidth}\raggedright
Requisito
\end{minipage} & \begin{minipage}[b]{\linewidth}\raggedright
Descrizione Operativa
\end{minipage} \\
\midrule\noalign{}
\endhead
\bottomrule\noalign{}
\endlastfoot
\textbf{\passthrough{\lstinline!SR-NET-01!}} & Gateway VPN WireGuard &
Allocazione di un'interfaccia kernel \passthrough{\lstinline!wg-k!} per
ciascun team \(k\), in ascolto sulla porta UDP \(51820 + k\) con
indirizzo gateway \passthrough{\lstinline!10.k.0.1/24!}. \\
\textbf{\passthrough{\lstinline!SR-NET-02!}} & Linux Bridge L2 Live &
Allocazione di un bridge Linux \passthrough{\lstinline!br-team-k!} per
ciascun team \(k\), con gateway
\passthrough{\lstinline!10.k.1.1/24!}. \\
\textbf{\passthrough{\lstinline!SR-NET-03!}} & Indirizzamento e
Credenziali VPN & Allocazione deterministica degli IP dei membri nella
subnet \passthrough{\lstinline!10.k.0.0/24!}. \\
\textbf{\passthrough{\lstinline!SR-NET-04!}} & Topologia Deterministica
& Configurazione statica degli indirizzi IP
\passthrough{\lstinline!10.k.1.x!} nei container in base al suffisso
(\passthrough{\lstinline!ip\_suffix!}) dichiarato nella specifica della
challenge. \\
\textbf{\passthrough{\lstinline!SR-NET-05!}} & TCP MSS Clamping &
Applicazione di regole \passthrough{\lstinline!nftables!} per forzare il
clamping del MSS, prevenendo pacchetti persi dovuti al tunnel
WireGuard. \\
\textbf{\passthrough{\lstinline!SR-NET-06!}} & Isolamento e Masquerading
NAT & Isolamento inter-team tramite il set
\passthrough{\lstinline!@team\_links!} e masquerading del traffico
outbound verso la rete esterna per l'accesso a Internet per giocatori,
container e rete di verifica. \\
\end{longtable}

\subsection{Motore di Verifica OVN
Sandbox}\label{motore-di-verifica-ovn-sandbox}

I requisiti del motore di verifica definiscono le garanzie di non
interferenza, trasparenza e sicurezza durante l'esecuzione degli exploit
e dei test di integrità.

\begin{longtable}[]{@{}
  >{\raggedright\arraybackslash}p{(\linewidth - 4\tabcolsep) * \real{0.1531}}
  >{\raggedright\arraybackslash}p{(\linewidth - 4\tabcolsep) * \real{0.2245}}
  >{\raggedright\arraybackslash}p{(\linewidth - 4\tabcolsep) * \real{0.6224}}@{}}
\toprule\noalign{}
\begin{minipage}[b]{\linewidth}\raggedright
Codice
\end{minipage} & \begin{minipage}[b]{\linewidth}\raggedright
Requisito
\end{minipage} & \begin{minipage}[b]{\linewidth}\raggedright
Descrizione Operativa
\end{minipage} \\
\midrule\noalign{}
\endhead
\bottomrule\noalign{}
\endlastfoot
\textbf{\passthrough{\lstinline!SR-VRF-01!}} & Switch Logico OVN &
Provisioning di uno switch virtuale OVN
(\passthrough{\lstinline!verify-<session>!}) senza esporre il gateway
all'host per evitare collisioni di rotte tra container. \\
\textbf{\passthrough{\lstinline!SR-VRF-02!}} & Fedeltà Topologica & I
cloni nella sandbox preservano l'interfaccia
\passthrough{\lstinline!eth0!} e la medesima configurazione IP
(\passthrough{\lstinline!10.k.1.x!}) specificata. \\
\textbf{\passthrough{\lstinline!SR-VRF-03!}} & Readiness Probes e
Timeout & Esecuzione preliminare dei controlli di prontezza sui servizi
con tentativi periodici, prima dell'esecuzione dei check di verifica. \\
\textbf{\passthrough{\lstinline!SR-VRF-04!}} & Confinamento Zero
Side-Channel & Isolamento del traffico di verifica sullo switch OVN per
impedire qualsiasi sniffing (es. \passthrough{\lstinline!tcpdump!}) da
parte dei concorrenti. \\
\textbf{\passthrough{\lstinline!SR-VRF-05!}} & Rilascio Automatico
Risorse & Deallocazione automatica e atomica di cloni, switch e porte
virtuali al termine della verifica o su eccezione. \\
\end{longtable}

\subsection{Control Plane, Concorrenza e
Persistenza}\label{control-plane-concorrenza-e-persistenza}

La tabella seguente specifica i vincoli di sincronizzazione, resilienza
e coerenza dello stato applicativo gestiti dal control plane di
BlueAgent.

\begin{longtable}[]{@{}
  >{\raggedright\arraybackslash}p{(\linewidth - 4\tabcolsep) * \real{0.1531}}
  >{\raggedright\arraybackslash}p{(\linewidth - 4\tabcolsep) * \real{0.2245}}
  >{\raggedright\arraybackslash}p{(\linewidth - 4\tabcolsep) * \real{0.6224}}@{}}
\toprule\noalign{}
\begin{minipage}[b]{\linewidth}\raggedright
Codice
\end{minipage} & \begin{minipage}[b]{\linewidth}\raggedright
Requisito
\end{minipage} & \begin{minipage}[b]{\linewidth}\raggedright
Descrizione Operativa
\end{minipage} \\
\midrule\noalign{}
\endhead
\bottomrule\noalign{}
\endlastfoot
\textbf{\passthrough{\lstinline!SR-CON-01!}} & Two-Tier Lock Registry &
Meccanismo di lock gerarchico a due livelli per prevenire race condition
sulle risorse. \\
\textbf{\passthrough{\lstinline!SR-CON-02!}} & Self-Healing e
Riconciliazione & Demone \passthrough{\lstinline!LifecycleManager!} per
la bonifica automatica degli switch e cloni orfani. \\
\textbf{\passthrough{\lstinline!SR-CON-03!}} & Healthcheck \& Telemetria
Live & Endpoint \passthrough{\lstinline!/health!} per ottenere lo stato
di salute del sistema. Raccolta telemetrica dello stato dei peer
WireGuard. \\
\textbf{\passthrough{\lstinline!SR-CON-04!}} & Persistenza Asincrona
SQLite & Memorizzazione dello stato di team, membri, container e
cataloghi su database SQLite. \\
\textbf{\passthrough{\lstinline!SR-CON-05!}} & Rollback Atomico su
Errore & Qualsiasi fallimento durante il provisioning di container o
interfacce di rete innesca la deallocazione delle risorse. \\
\textbf{\passthrough{\lstinline!SR-CON-06!}} & Idempotenza Operativa &
Tutte le richieste di creazione e configurazione per risorse già
esistenti non generano errori applicativi. \\
\textbf{\passthrough{\lstinline!SR-CON-07!}} & Autenticazione
Centralizzata & Validazione obbligatoria dell'accesso a tutti gli
endpoint tramite token segreto condiviso. \\
\end{longtable}

\subsection{Pipeline Immagini
Distrobuilder}\label{pipeline-immagini-distrobuilder}

I requisiti della pipeline stabiliscono la creazione riproducibile delle
immagini per le sfide, il loro tracciamento e il caching dei binari.

\begin{longtable}[]{@{}
  >{\raggedright\arraybackslash}p{(\linewidth - 4\tabcolsep) * \real{0.1531}}
  >{\raggedright\arraybackslash}p{(\linewidth - 4\tabcolsep) * \real{0.2245}}
  >{\raggedright\arraybackslash}p{(\linewidth - 4\tabcolsep) * \real{0.6224}}@{}}
\toprule\noalign{}
\begin{minipage}[b]{\linewidth}\raggedright
Codice
\end{minipage} & \begin{minipage}[b]{\linewidth}\raggedright
Requisito
\end{minipage} & \begin{minipage}[b]{\linewidth}\raggedright
Descrizione Operativa
\end{minipage} \\
\midrule\noalign{}
\endhead
\bottomrule\noalign{}
\endlastfoot
\textbf{\passthrough{\lstinline!SR-IMG-01!}} & Definizione IaC
Distrobuilder & Creazione dichiarativa di immagini per container di
sistema a partire da manifesti YAML versionabili. \\
\textbf{\passthrough{\lstinline!SR-IMG-02!}} & Caching e Identificazione
SHA-256 & Tracciamento crittografico della specifica YAML tramite hash
SHA-256 per identificare modifiche su immagini già compilate nel
database. \\
\textbf{\passthrough{\lstinline!SR-IMG-03!}} & Compilazione Asincrona &
Esecuzione delle fasi di build e import Incus in background tramite task
asincrono con tracciamento dei log. \\
\end{longtable}
